\section{SECCIÓN 3. TAREAS DE SQA}
El SQA se encargará de evaluar la calidad del sistema tanto en la capa de presentación como en los endpoints del sistema, con mediciones de rendimiento, eficiencia, seguridad, efectividad y corrección.

\subsection{3.1 TAREA: REVISIÓN DE PRODUCTOS DE SOFTWARE}
Dentro del sistema evaluado, se encontró poca documentación para analizar; apenas se definen las funciones principales, las tecnologías utilizadas y la configuración inicial del sistema. No tiene pruebas definidas para comprobar la total funcionalidad del sistema, por lo que la mayoría de las pruebas y evaluaciones se realizarán a partir de las pruebas adaptadas al software entregado a través del plan generado para dichas pruebas.

\subsection{3.2 TAREA: REVISIÓN DE HERRAMIENTAS DE SOFTWARE}
Se definirán las herramientas utilizadas para las mediciones, tanto en qué parte del proceso se utilizarán como para qué tipos de prueba se usarán. Todas estas pruebas serán nuevas en el proyecto al no tener definido inicialmente antes del análisis. Cada una está pensada para el testeo en el ambiente en el que el sistema fue desarrollado, como son Java y TypeScript, que definen el servicio web.

\subsection{3.3 TAREA: EVALUAR PROCESO DE REVISIÓN DE PRODUCTOS DE SOFTWARE}
Esta tarea de SQA asegura que los procesos de revisión de calidad estén establecidos para todos los productos de software, que pueden incluir representaciones de información distintas a los documentos tradicionales en papel, y que estos productos hayan pasado por evaluación, pruebas y acciones correctivas según lo requiera el estándar.

SQA verificará que los productos de software listos para revisión sean revisados, que los resultados se reporten y que los problemas reportados se resuelvan de acuerdo con la referencia (e).

Los resultados de esta tarea se documentarán utilizando el Formulario de Auditoría de Procesos descrito en la Sección 7 y se proporcionarán a la gestión del proyecto. La recomendación de SQA para la acción correctiva requiere la disposición de la gestión del proyecto y se procesará de acuerdo con la guía en la Sección 7.

\subsection{3.4 TAREA: EVALUAR PROCESOS DE PLANIFICACIÓN, MONITOREO Y SUPERVISIÓN DE PROYECTOS}
La planificación, seguimiento y supervisión del proyecto implica que la gestión del proyecto desarrolle y documente planes para el Desarrollo de Software, Pruebas de CI y Sistemas, Instalación de Software y Transición de Software. La Sección 1.6 enumera los planes del programa/proyecto. Para los documentos del proyecto que se desarrollarán, SQA ayudará a identificar las pautas, estándares o Descripciones de Elementos de Datos (DIDs) apropiados, y ayudará con la adaptación de esas pautas, estándares o DIDs para satisfacer las necesidades del proyecto.

SQA evaluará que el proyecto realice las actividades relevantes establecidas en los planes del Programa y del Proyecto. Para verificar que estas actividades se realicen según lo planificado, SQA auditará los procesos que definen la actividad y utilizará la referencia (e) u otro documento de planificación como medida de cumplimiento.

SQA utilizará las listas de verificación de auditoría en las Figuras B-1 y B-2 como guías para realizar la evaluación.

Los resultados de esta tarea se documentarán utilizando el Formulario de Auditoría de Procesos descrito en la Sección 7. Cualquier cambio recomendado a esos planes requerirá actualización y aprobación por parte de la gestión del proyecto de acuerdo con el procedimiento de gestión de configuración descrito en el Plan SCM del proyecto.

\subsection{3.5 TAREA: EVALUAR PROCESOS DE IMPLEMENTACIÓN, PRUEBAS UNITARIAS E INTEGRACIÓN}
La implementación o codificación del software es el punto en el ciclo de vida del desarrollo en el que el diseño finalmente se implementa. El proceso incluye pruebas unitarias del código de software, integración de componentes y pruebas de calificación de CI.

En la fase de integración y prueba del ciclo de vida del desarrollo, el enfoque de las pruebas cambia de la corrección de un componente individual al funcionamiento adecuado de las interfaces entre componentes, el flujo de información a través del sistema y el cumplimiento de los requisitos del sistema.

Las actividades de SQA se enumeran a continuación:
\begin{enumerate}
    \item Verificar que el proceso de codificación, las revisiones de código asociadas y las pruebas unitarias de software se realicen de acuerdo con los estándares y procedimientos establecidos por el proyecto y descritos en la Sección 5.
    \item Verificar que los elementos de acción resultantes de las revisiones del código se resuelvan de acuerdo con estos estándares y procedimientos.
    \item Verificar que las actividades de prueba de software estén identificadas, que los entornos de prueba estén definidos y que se hayan diseñado pautas para las pruebas.
    \item Verificar que el proceso de integración de software, las actividades de prueba de integración de software y las actividades de prueba de rendimiento del software se realicen de acuerdo con el plan de pruebas y los estándares del proyecto.
    \item Verificar que se presencien las pruebas de integración de software y las pruebas de rendimiento del software para verificar que se sigan los procedimientos de prueba aprobados, que se mantengan registros precisos de los resultados de las pruebas, que todas las discrepancias descubiertas durante las pruebas se informen correctamente y que se completen los informes de prueba asociados.
    \item Verificar que las discrepancias descubiertas durante las pruebas de integración y rendimiento del software se identifiquen, analicen, documenten y corrijan; que las pruebas unitarias y de integración del software se vuelvan a ejecutar según sea necesario para validar las correcciones realizadas al código.
    \item Verificar que la responsabilidad de las pruebas y de la presentación de informes sobre los resultados se haya asignado a un elemento organizativo específico.
    \item Revisar el Plan de Pruebas de Software y las Descripciones de Pruebas de Software para verificar el cumplimiento de los requisitos y estándares.
    \item Monitorear las actividades de prueba, presenciar pruebas y certificar los resultados de las pruebas.
    \item Verificar que se hayan establecido requisitos para la certificación o calibración de todo el software o hardware de soporte utilizado durante las pruebas.
\end{enumerate}

SQA puede utilizar las listas de verificación de auditoría en las Figuras B-7, B-8 y B-9 como guías para realizar estas evaluaciones.

Los resultados de esta tarea se documentarán utilizando el Formulario de Auditoría de Procesos descrito en la Sección 7 y se proporcionarán a la gestión del proyecto. La recomendación de SQA para la acción correctiva requiere la disposición de la gestión del proyecto y se procesará de acuerdo con la guía en la Sección 7.

\subsection{3.6 TAREA: EVALUAR PROCESOS DE CONTROL DE CÓDIGO}
El control de código incluye los métodos e instalaciones utilizados para mantener, almacenar, asegurar y documentar las versiones controladas del software identificado durante todas las fases del ciclo de vida del software. El proyecto utiliza GitHub como repositorio remoto para el control de versiones del código fuente.

Las actividades de SQA se enumeran a continuación:
\begin{enumerate}
    \item Verificar que el repositorio de código esté establecido y que existan procedimientos para gobernar su operación (ramas, pull requests, merge).
    \item Verificar que la documentación y los materiales del programa informático estén aprobados y colocados bajo control de gestión de configuración en el repositorio.
    \item Verificar el establecimiento de procedimientos de liberación formal para la documentación y versiones de software aprobadas.
    \item Verificar que los controles del repositorio prevengan cambios no autorizados en el software controlado (protección de ramas, revisiones obligatorias) y verificar la incorporación de todos los cambios aprobados.
\end{enumerate}

SQA puede utilizar la lista de verificación de auditoría en la Figura B-18 como guía para realizar esta evaluación.

Los resultados de esta tarea se documentarán utilizando el Formulario de Auditoría de Procesos descrito en la Sección 7 y se proporcionarán a la gestión del proyecto. La recomendación de SQA para la acción correctiva requiere la disposición de la gestión del proyecto y se procesará de acuerdo con la guía en la Sección 7.

Los informes de SQA, junto con los registros de acción correctiva, los informes de pruebas de software y los registros de evaluación de productos de software, constituyen la evidencia requerida para certificar que el software desempeña las funciones previstas.

\subsection{3.7 TAREA: VERIFICAR REVISIONES DE PROCESOS Y AUDITORÍAS}
El tipo y alcance de las revisiones técnicas depende en gran medida del tamaño, alcance, riesgo y criticidad del proyecto de software. Para este proyecto académico, se definen las siguientes revisiones y auditorías adaptadas al alcance del sistema.

\subsubsection{ Tarea: Verificar Revisiones Técnicas}
Un componente principal de la ingeniería de calidad en el software es la realización de revisiones técnicas de los productos de software, tanto entregables como no entregables. Los participantes de una revisión técnica deben incluir personas con conocimiento técnico de los productos de software a revisar. El propósito de la revisión técnica será centrarse en los productos de software en progreso y finales, en lugar de los materiales generados especialmente para la revisión.

Para este proyecto se definen las siguientes revisiones técnicas:

\begin{enumerate}
    \item \underline{Revisión por Pares}: el objetivo es revisar el código fuente, la documentación y los productos de software en progreso para identificar defectos, verificar el cumplimiento de estándares y obtener retroalimentación técnica. Se realizará antes de fusionar cambios en la rama principal.
    \item \underline{Revisión de Diseño}: el objetivo es evaluar la arquitectura del sistema (backend Spring Boot, frontend React, persistencia PostgreSQL), las interfaces entre componentes y la adecuación de la solución técnica a los requisitos de seguridad y funcionalidad.
\end{enumerate}

Las revisiones técnicas se llevarán a cabo para revisar los productos de software en evolución, demostrar las soluciones técnicas propuestas y proporcionar información y obtener retroalimentación sobre el esfuerzo técnico. El resultado de una revisión técnica se enumera a continuación:
\begin{enumerate}
    \item Identificar y resolver problemas técnicos.
    \item Revisar el estado del proyecto, específicamente los riesgos a corto y largo plazo en cuanto a aspectos técnicos, costos y plazos.
    \item Llegar a estrategias de mitigación acordadas para los riesgos identificados, dentro de la autoridad de los presentes.
    \item Identificar riesgos y problemas que se plantearán en las revisiones conjuntas de gestión.
    \item Verificar la comunicación continua entre los miembros del equipo del proyecto.
\end{enumerate}

Los criterios de entrada para una revisión técnica requerirán que un elemento a revisar se distribuya al grupo antes de la reunión de revisión. Además, se asignará un secretario para registrar cualquier problema que requiera resolución, indicando el responsable del elemento de acción y la fecha de vencimiento, y las decisiones tomadas dentro de la autoridad de los participantes de la revisión técnica.

Se recopilan varias mediciones como parte de las revisiones técnicas para ayudar a determinar la efectividad del proceso de revisión en sí, así como de los pasos del proceso que se utilizan para producir el elemento que se revisa. Estas mediciones, reportadas al gerente del proyecto, incluirán la cantidad de tiempo dedicado por cada persona involucrada en la revisión, incluida la preparación para la revisión.

\subsubsection{ Tarea: Verificar Revisiones de Gestión}
La revisión periódica de la gestión de SQA del estado, progreso, problemas y riesgos del proyecto de software proporcionará una evaluación independiente de las actividades del proyecto. SQA proporcionará la siguiente información a la gestión:
\begin{enumerate}
    \item \textbf{Cumplimiento:} Identificación del nivel de cumplimiento del proyecto con los procesos organizacionales y del proyecto establecidos.
    \item \textbf{Áreas problemáticas:} identificación de áreas problemáticas potenciales o reales del proyecto basadas en el análisis de los resultados de las revisiones técnicas.
    \item \textbf{Riesgos:} identificación de riesgos basados en la participación y evaluación del progreso del proyecto y las áreas problemáticas.
\end{enumerate}

Debido a que la función SQA es integral para el éxito del proyecto, SQA comunicará libremente sus resultados a la dirección del proyecto y al equipo del proyecto. El método para informar el cumplimiento, las áreas problemáticas y los riesgos se comunicará en un informe documentado o memorando. Se dará seguimiento a los problemas de cumplimiento, áreas problemáticas y riesgos y se rastrearán hasta su cierre.

\subsubsection{ Tarea: Conducir Auditorías de Procesos}
Los procesos de desarrollo de software se auditan de acuerdo con las tareas especificadas en esta Sección y se realizan de acuerdo con el cronograma de desarrollo de software. SQA utilizará las listas de verificación de auditoría del Apéndice B como guías para realizar las evaluaciones.

Los resultados de esta tarea se documentarán utilizando el Formulario de Auditoría de Procesos descrito en la Sección 7 y se proporcionarán a la gestión del proyecto.

\subsubsection{ Tarea: Realizar Auditorías de Configuración}
\textbf{Auditoría de Configuración Funcional (FCA).} La Auditoría de Configuración Funcional se lleva a cabo antes de la entrega del software para comparar el software construido con los requisitos de software establecidos. El propósito es asegurar que el código aborda todos, y solo, los requisitos documentados y las capacidades funcionales establecidas. SQA participará como miembro del equipo de FCA. SQA ayudará en la preparación de los hallazgos de la FCA. Cualquier seguimiento de los hallazgos reportados de la FCA será monitoreado y rastreado hasta su cierre.

\textbf{Auditoría de Configuración Física (PCA).} La Auditoría de Configuración Física se lleva a cabo para verificar que el software y su documentación son internamente consistentes y están listos para la entrega. El propósito es asegurar que la documentación que se va a entregar sea consistente y describa correctamente el código. SQA participará como miembro del equipo de PCA. SQA ayudará en la preparación de los hallazgos de la PCA. Cualquier seguimiento de los hallazgos reportados de la PCA será monitoreado y rastreado hasta su cierre.

\subsection{3.8 RESPONSABILIDADES}
La responsabilidad última de la calidad del software y la documentación asociada del proyecto recae en el Gerente de Proyecto de Software del proyecto. El Gerente de SQA implementará los procedimientos de SQA definidos en este plan.

SQA deriva su autoridad del Gerente del Proyecto. SQA monitoreará las actividades del personal del proyecto y revisará los productos para verificar el cumplimiento de los estándares y procedimientos aplicables. Los resultados del monitoreo y análisis de SQA, junto con las recomendaciones de SQA para acciones correctivas, se reportarán al Gerente de Proyecto de Software. Todos los documentos y software aprobados por el Gerente de Proyecto de Software para su liberación deberán haber sido revisados y aprobados por SQA. La Tabla 3-1 es una matriz de responsabilidades para las tareas identificadas en esta Sección.

\begin{table}[H]
\centering
\caption{Matriz de Responsabilidades por Proceso}
\label{tab:responsibility}
\begin{tabular}{|p{5cm}|p{1.5cm}|p{1.5cm}|p{1.5cm}|p{1.3cm}|p{1.3cm}|p{1.5cm}|}
\hline
\textbf{Proceso} & \textbf{Ger. SQA} & \textbf{Ger. Prog.} & \textbf{Ger. Proy.} & \textbf{Ing. Sist.} & \textbf{Des. SW} & \textbf{Probador} \\
\hline
Desarrollar/Documentar Plan SQA & X & & X & & & \\
\hline
Revisar Plan SQA & X & X & X & X & X & X \\
\hline
Aprobar Plan SQA & X & X & X & & & \\
\hline
Revisar productos de software & X & X & X & X & X & X \\
\hline
Aprobar producto & & X & X & & & \\
\hline
Evaluar herramientas de software & X & & & & X & X \\
\hline
Desarrollar/Documentar planes del proyecto & & X & X & & & \\
\hline
Revisar planes del proyecto & X & X & X & X & X & X \\
\hline
Evaluar proceso de planificación & X & & & & & \\
\hline
Desarrollar/Corregir código & & & & & X & \\
\hline
Revisar código & X & & & & X & X \\
\hline
Prueba unitaria & & & & & X & X \\
\hline
Integrar software & & & & & X & \\
\hline
Probar software integrado & & & & & & X \\
\hline
Prueba de rendimiento & & & & & & X \\
\hline
Evaluar/Reportar proceso de implementación y pruebas & X & & & & & \\
\hline
Control de código (GitHub) & X & & & & X & \\
\hline
Revisar requisitos del sistema & X & X & X & X & X & X \\
\hline
Revisar requisitos de software & X & X & X & X & X & X \\
\hline
Revisar diseño de software & X & X & X & X & X & X \\
\hline
Auditorías de procesos & X & & & & & \\
\hline
Auditorías de configuración & X & & & X & X & X \\
\hline
Seguir proceso de acción correctiva & X & X & X & X & X & X \\
\hline
Resolver hallazgos de auditoría & & X & X & & & \\
\hline
Asistir a auditorías SQA & X & & & X & X & X \\
\hline
Nombrar auditor SQA independiente & & X & & & & \\
\hline
\end{tabular}
\end{table}

\subsection{3.9 CRONOGRAMA}
Los cronogramas de SQA están estrechamente coordinados con el cronograma de desarrollo de software. El proyecto se ejecutó en un período de dos semanas con las siguientes fases:

\begin{table}[H]
\centering
\caption{Cronograma del Proyecto}
\label{tab:schedule}
\begin{tabular}{|p{2.5cm}|p{4cm}|p{6.5cm}|}
\hline
\textbf{FASE} & \textbf{ACTIVIDADES SQA} & \textbf{PRODUCTOS} \\
\hline
Semana 1: Revisión y Análisis & Revisión del sistema existente, identificación de herramientas SQA, definición del plan de pruebas, auditoría de procesos iniciales & SQAP, Plan de Pruebas (STP), definición de CP01--CP38, matriz de trazabilidad \\
\hline
Semana 2: Pruebas y Documentación & Ejecución de casos de prueba (CP01--CP38), generación de informes (JaCoCo, SonarQube, ZAP, Lighthouse, Cypress), registro de hallazgos y acciones correctivas, elaboración de documentación SQA & Informe de Pruebas (STR), cobertura JaCoCo, hallazgos SEC-01/XSS-01/CFG-01/BLD-01/CAL-01, RESULTADOS\_CP30\_CP39 \\
\hline
\end{tabular}
\end{table}

Las auditorías de procesos se realizarán al comienzo de cada nueva fase de desarrollo para verificar que los procesos apropiados estén correctamente implementados según lo definido en los documentos de planificación. Además, se realizarán verificaciones puntuales (auditorías no programadas) durante cada fase de desarrollo para verificar que se sigan los procesos y los procedimientos de escritorio. Al finalizar una fase de desarrollo de software, SQA revisará e informará si se han completado todos los pasos necesarios para la transición a la siguiente fase.
